\chapter{Advanced GM Techniques}
łabel{ch:gm−advanced}
∈dex{advanced techniques}

This chapter assumes you have mastered Fate's Edge core rules. It offers tools for GMs running \textbf{experienced groups and higher−tier campaigns (Tier III+)}. If you are new to the game, start with earlier chapters.

\section{Advanced Story Beat Management}
∈dex{Story Beats!advanced}

Story Beats (SB) are your primary tool for escalating tension. At high tiers, you will have more SB to spend −− but also more narrative threads to manage.

\begin{tcolorbox}[
  colback=blue!5!white,
  colframe=blue!75!black,
  title=You Are Not Limited by SB,
  fonttitle=\bfseries
]
∈dex{Story Beats!as pacing tool}
\textbf{SB are a pacing tool, not a budget.} If the story needs a complication, introduce it. SB just tells you when players \textit{want} escalation (by rolling 1s) and gives you a reservoir to draw from. Do not hoard SB; spend them freely to keep scenes dynamic.

\medskip
\noindent\textit{Example:} The party is interrogating a prisoner. Even without a roll, you sense the scene dragging. Spend 1 SB: ``The prisoner's eyes dart to the window −− a signal. Someone is listening.'' No roll, no budget −− just pacing.
\end{tcolorbox}

\subsection*{Spending SB Across Multiple Threads}
Instead of pouring all SB into one complication, distribute them across several active timers or fronts (see \ref{sec:core−resolution−cycle}). A scene might have:
\begin{itemize}
  ℹtem \emph{Alarm Raised [4]} −− advances from noise or failed stealth.
  ℹtem \emph{Rival Closing In [6]} −− advances from delays or loud magic.
  ℹtem \emph{Patron Attention [8]} −− advances each time a Rite is invoked.
\end{itemize}
Each Miss or Partial can tick one or two small timers instead of a single big setback. This keeps pressure distributed and gives players multiple problems to solve.

\subsection*{Using SB to Advance Campaign Fronts}
When a scene directly relates to a long−term front (e.g., a war, a spreading curse), spend SB to advance that front's timer −− even if the scene is otherwise a ``success''. Example: The party defeats a cultist cell (clean success), but you spend 2 SB to tick \emph{Cult Influence [6]} −− their network still grows elsewhere.

\subsection*{Banked SB from Player Complications}
Characters may start sessions with banked SB from unresolved complications (see). These count toward scene spending limits but should be spent on \textbf{thematic complications tied to that character's background or Patron}, not generic setbacks. This makes personal storylines matter.

\subsection*{Example: Cascade Spend}
A single Miss generates 2 SB. The GM:
\begin{enumerate}
  ℹtem Spends 1 SB to advance \emph{Alarm Raised [4]} by 1 (now 2/4 −− guards are converging).
  ℹtem Spends 1 SB to advance \emph{Rival Closing In [6]} by 1 (now 3/6 −− rival scouts are spotted).
  ℹtem Both timers are now half full −− the scene's tension doubles without a single ``big'' complication.
\end{enumerate}

\section{Mastering the Deck of Consequences}
∈dex{Deck of Consequences!advanced}

The deck is not just a randomizer; it is a \textbf{thematic engine}. Use it to keep complications fresh and unpredictable, even at high tiers.

\subsection*{Shortcut Table: Suit + Rank → Typical Complication}
łabel{subsec:deck−shortcut}
Use this table for fast interpretation at the table.

\begin{table}[htbp]
¢ering
\small
\begin{tabularx}{\textwidth}{lXX}
⊤rule
\textbf{Suit} & \textbf{Low (2−−5)} & \textbf{High (J−−A)} \\
∣rule
Hearts (Emotion/Social) & Awkward social moment, minor suspicion & Betrayal, scandal, broken trust \\
Spades (Physical/Harm) & Bruise, small gear damage, stumble & Serious injury, structural collapse \\
Clubs (Resource/Cost) & Minor loss (rations, a few coins) & Major asset destroyed, debt called \\
Diamonds (Magic/Spiritual) & Omen, eerie noise, minor curse & Patron demand, reality bend \\
⊥tomrule
\end{tabularx}
∩tion{Shortcut table: Suit + Rank → Typical Complication}
łabel{tab:deck−shortcut}
\end{table}

\subsection*{Combining Draws (Suit Synthesis)}
Draw up to 3 cards per SB spend (minimum 1, maximum 3). Synthesize a single twist from the combination of suits, not just the highest rank.

\begin{itemize}
  ℹtem \textbf{Hearts + Spades:} Emotional betrayal that leads to physical harm (``Your ally flees, and the portcullis drops on you.'')
  ℹtem \textbf{Clubs + Diamonds:} Resource drain with magical backlash (``The ritual consumes your healing herbs and leaves a tainted scar.'')
  ℹtem \textbf{Hearts + Clubs + Spades:} Social pressure plus resource loss plus injury −− a crowd riots and loots your gear.
\end{itemize}

\subsection*{Using the Deck for Non−Linear Scenes}
Draw a card to inspire flashbacks, montages, or off−screen events. Example: During downtime, draw a Clubs card −− ``A fire at your safehouse.'' Instead of playing it, you narrate a short scene of the party arriving to find the aftermath. The deck becomes a prompt generator, not merely a roll−aftermath tool.

\subsection*{Custom Suit Tables}
For a specific campaign or region, replace the default suit meanings with custom ones. Example −− Shadow Court campaign:
\begin{itemize}
  ℹtem \textbf{Hearts:} Vows and obligations.
  ℹtem \textbf{Spades:} Physical traps and barriers.
  ℹtem \textbf{Clubs:} Political pressure and blackmail.
  ℹtem \textbf{Diamonds:} Oaths and geasa.
\end{itemize}
This reflavors the deck without changing its structure.

\section{Adjudicating Thresholds}
łabel{sec:thresholds−advanced}
∈dex{thresholds!adjudication}

Thresholds −− doorways, crossroads, riverbanks, the space between waking and sleeping −− are where the world's memory is thinnest. When a player announces ``I cross the threshold,'' you need a consistent way to respond.

\subsection*{Common Threshold Types and Their Effects}
Use this table as a reference. The effects are narrative nudges, not rigid rules.

\begin{table}[htbp]
¢ering
\small
\begin{tabularx}{\textwidth}{lXX}
⊤rule
\textbf{Threshold} & \textbf{Typical Effect} & \textbf{SB Spend Opportunity} \\
∣rule
Doorway (unwarded) & Position worsens one step (attention from spirits) & 1 SB: door leads to wrong room \\
Doorway (warded) & Resolve test (DV 3) or mark Fatigue & 2 SB: ward triggers (Harm 1) \\
Crossroads & GM asks a ``choice'' question; answer becomes true & 3 SB: a stranger appears with a bargain \\
Riverbank & Roll Spirit + Survival; success gives omen & 2 SB: spirit from the water surfaces \\
Bridge & The bridge ``remembers'' −− ask a past event & 3 SB: bridge's toll is a memory \\
Sleep/waking edge & Player narrates a dream; GM gains 1 SB & 4 SB: dream becomes real for a scene \\
⊥tomrule
\end{tabularx}
∩tion{Common threshold types and their effects}
łabel{tab:thresholds}
\end{table}

\subsection*{Threshold Timers}
For a location that is heavily liminal (a haunted manor, a fey crossing), start a \textbf{Threshold Instability [4−−6]} timer. Tick it each time a player crosses a threshold without proper respect (no offering, no announced name). When the timer fills, the threshold inverts: doors open to the wrong places, the dead walk, or the party is lost in a mirror realm.

\subsection*{Player−Driven Thresholds}
When a player says ``I act at the threshold,'' encourage them to be specific: ``I open the door with my left hand'' or ``I step over the river stone.'' Use that fiction to set Position. A respectful crossing (removing boots, pouring a libation) may start Controlled; a brash crossing starts Desperate.

\section{Faction Grand Strategy}
∈dex{factions!advanced}

At high tiers, factions act as ``campaign engines.'' Build them with \textbf{faction timers} that advance whether the players intervene or not.

\subsection*{Faction Timer Examples}
\begin{itemize}
  ℹtem \textbf{War [8]} −− Advances when battles occur or supply lines are cut.
  ℹtem \textbf{Influence [6]} −− Advances when allies are won or propaganda spreads.
  ℹtem \textbf{Intrigue [8]} −− Advances when secrets are uncovered or betrayals happen.
\end{itemize}
These timers are \textbf{visible to players} −− they see the war brewing, even if they cannot stop it alone.

\subsection*{Player Actions and Faction Timers}
When players act, they can tick a faction timer forward or backward. Example: Sabotaging a siege engine ticks \emph{War −−1} (enemy's timer) and ticks \emph{Their Own Mandate +1} (see \ref{subsec:mandate}). Document these changes openly. The world becomes a reactive board.

\subsection*{Running a ``Faction Turn'' Between Sessions}
łabel{subsec:faction−turn}
For a highly political campaign, between sessions you can:
\begin{itemize}
  ℹtem Ask players for their faction's orders (e.g., ``Send spies to the capital'').
  ℹtem Roll a single die per order (Pass/Fail/Partial) using the faction's relevant Asset (e.g., Spy Ring).
  ℹtem Advance faction timers accordingly.
  ℹtem Start the next session with a brief ``news from the front'' segment.
\end{itemize}
This makes factions feel alive without heavy bookkeeping.

\subsection*{Example: Faction Turn with Rivals}
The rival party \textbf{The Ashen Cord} has their own faction timer: \emph{Cord's Agenda [6]}. During the faction turn:
\begin{itemize}
  ℹtem The GM rolls for the Cord's order: ``Secure the Sunstone fragment.''
  ℹtem The roll is a Partial. The Cord finds the fragment but leaves evidence (tracks, a witness).
  ℹtem The GM ticks \emph{Cord's Agenda} by 1 and notes: ``A merchant in Silkstrand saw a scarred woman matching Vesper's description.''
  ℹtem At the next session, the players hear the rumor. They now have a choice: pursue the Cord or stay on their own path.
\end{itemize}

\section{Advanced Timer Design}
∈dex{timers!advanced}

Timers become more powerful when they interact. Use the following patterns to create tension that players must actively manage.

\subsection*{Interacting Timers}
Two timers that affect each other. Example: \emph{Ritual Progress [6]} and \emph{Cult Interference [4]}. Ticking one may automatically tick the other, or players can choose which to slow.

\subsection*{Cascading Timers}
When one timer fills, it creates another. Example: \emph{Security Alert [4]} fills → \emph{Reinforcements Arrive [6]} appears. This models escalating threats without requiring the GM to track everything from the start.

\subsection*{Hidden Consequences Behind Success}
Even a clean success can tick a long−term timer. Example: The party assassinates a cult leader (clean success). Off−screen, \emph{Cult of the Martyr [8]} ticks +2 −− the death radicalizes the remaining followers.

\section{Designing High−Tier Custom Content}
∈dex{custom content!high tier}

At Tiers IV−−V, characters can reshape nations and realities. Custom content should grant \textbf{narrative permissions}, not merely bigger numbers.

\subsection*{Epic Talents (12+ XP)}
An epic talent might:
\begin{itemize}
  ℹtem Allow rewriting one line of a Patron's Rite (changing its cost or effect).
  ℹtem Grant a once−per−session ``truth edit'' −− the GM must accept one factual change to the scene.
  ℹtem Let a character invoke a Patron's name to retroactively avoid a TPK (once per campaign, with heavy cost).
\end{itemize}
These are not combat bonuses −− they change what kind of story is being told.

\subsection*{Assets for Ruling Nations}
A Major Asset (12 XP) might be a kingdom. Its free use could be ``order a levy'' or ``declare a holiday''. Its scene surge (1 Boon) could be ``call in a noble favor'' or ``divert a river''. Such assets come with obvious upkeep (taxes, succession, rebellions) −− and they can be threatened or lost.

\subsection*{Adventures for Demigods}
At the highest tiers, stakes should threaten \textbf{concepts}, not just lives. Examples:
\begin{itemize}
  ℹtem A curse that erases the memory of a city.
  ℹtem A war that will redraw the map for centuries.
  ℹtem A Patron who will change the fundamental laws of magic.
\end{itemize}
Victory may not be guaranteed −− but even failure should leave a permanent mark on the world.

\section{Running Mythic Scenarios (Tiers IV−−V)}
∈dex{mythic scenarios}

At the apex of power, the game changes. Use these techniques to keep play grounded in choices, not points.

\subsection*{Handling Reality−Altering Magic}
When a player attempts something cosmic (e.g., ``I wish the volcano had never erupted''), do not say no. Instead, offer a \textbf{compromise}:
\begin{itemize}
  ℹtem ``You can do it, but it will cost you your connection to this place −− you will never return.''
  ℹtem ``It works, but for the remainder of the session, the world is unstable −− every roll generates an extra SB.''
  ℹtem ``You succeed, but your Patron demands a permanent sacrifice −− one of your Talents is erased.''
\end{itemize}
The effect happens, but the price is felt.

\subsection*{Player Agency at Cosmic Scale}
\begin{tcolorbox}[
  colback=green!5!white,
  colframe=green!75!black,
  title=At Tier V, Every Choice May Erase an Alternative Timeline,
  fonttitle=\bfseries
]
∈dex{mythic scenarios!sacrifice}
Embrace this. Let players narrate what they are saving (or sacrificing) when they act.

\medskip
\noindent\textit{Example:} ``I stop the ritual −− but I let the rival escape.'' That rival becomes a future problem. ``I save the city −− but I lose my connection to my Patron.'' Now the player has a new arc: how to regain favor or find a new source of power.

\medskip
\noindent\textbf{Let the world break.} If the climax involves a collapsing reality, the battlefield itself can be a timer. Example: \emph{Reality Fracture [8]}. Each time a major event occurs (a spell, a death, an oath), tick the timer. When it fills, the scene shifts −− time loops, gravity reverses, or the characters meet themselves from another timeline.
\end{tcolorbox}

\subsection*{Let the World Break (Timer Example)}
∈dex{timers!reality fracture}
Set a visible timer: \textbf{Reality Fracture [8]}. Tick it when:
\begin{itemize}
  ℹtem A character invokes a Patron's name in desperation.
  ℹtem A major magical effect is attempted.
  ℹtem A character makes a life−or−death choice.
\end{itemize}
At 4 segments, strange echoes appear (ghostly duplicates, time slips). At 8 segments, the scene transforms: gravity inverts, doors lead to wrong locations, or the party faces an alternate version of themselves.

\section{Managing Cross−Patron and Invoker Conflicts}
∈dex{Patrons!cross−Patron conflicts}
∈dex{Invoker!Symbol friction}

When a character serves multiple Patrons (e.g., an Invoker with several Symbols, or a Runekeeper who also carries a Symbol), their obligations can conflict.

\subsection*{Signs of Cross−Resonance}
\begin{itemize}
  ℹtem A Symbol grows warm when another is used.
  ℹtem Dreams become fragmented −− two Patrons arguing in metaphor.
  ℹtem Rites from one Patron cost +1 Obligation because the other Patron is ``watching.''
\end{itemize}

\subsection*{Resolving Conflicts Mechanically}
If a character invokes two conflicting Patrons in the same scene, the GM may:
\begin{itemize}
  ℹtem Increase the DV of the second Rite by 1 (the Weave resists).
  ℹtem Tick a \emph{Patron Jealousy [4]} timer. When it fills, one Patron demands the character renounce the other (or pay a heavy price).
  ℹtem Offer a Devil's Bargain: ``You can use both, but mark 2 Fatigue and the GM gains 2 SB.''
\end{itemize}

\subsection*{Narrative Resolution}
Encourage players to roleplay the tension. A character who serves both the Traveler and Malachai might struggle with whether to take a shortcut (Traveler) that feeds a cursed addiction (Malachai). Use these moments to advance character arcs, not just punish mechanics.

\section{Player−Managed Modules at High Tiers}
∈dex{player−managed modules!high tier}

At high tiers, players track more resources: multiple Obligation segments per Patron, Corruption timers, Asset conditions, and often a Cord or Hollowed. Remind the table of the following responsibilities (see also Section 2.6 of the \textit{GM Guide}):

\begin{itemize}
  ℹtem \textbf{Players declare thresholds.} ``My Obligation just hit capacity'' triggers an immediate effect. The GM should not have to ask.
  ℹtem \textbf{Players maintain their own Asset states.} If a Safehouse is neglected, the player should note it and inform the GM when it is used.
  ℹtem \textbf{Players track their own Corruption Blooms.} When a Cantor's Corruption Timer fills, the player chooses (or rolls) the new trait and enforces its roleplay cost.
\end{itemize}

The GM's job is to \textbf{spot−check} −− occasionally ask ``What's your current Obligation to Malachai?'' −− and to \textbf{enforce consequences} when a threshold is crossed. Never hunt for a player's misplaced number. If they forget a resource, assume it is in their favor (they have it) but the story may suffer a small twist.

\section{Advanced Complex Scenarios}

\subsection*{Heists: Player−Designed Plans}
Instead of planning the heist for them, ask each player: ``How does your character contribute to the plan?'' Build a timer (e.g., \emph{Heist Progress [8]}). Each player makes a roll using their chosen method (Stealth, Deception, Arcana, etc.). Success ticks the timer; failure adds complications (alarm, guard alerted). The GM spends SB from misses to introduce obstacles from the security system.

\subsection*{Mass Combat: Unit Timers}
Instead of rolling for every soldier, treat opposing armies as a single \emph{War Timer [8]}. Player actions (command, sabotage, rally) tick the timer. When it fills, that side breaks or routs. Personal combat can happen in the same scene −− but the mass conflict is abstracted to a single visible tracker.

\subsection*{Political Intrigue: Reputation as a Resource}
Track a \emph{Reputation Timer [6]} for the party in a given court. Social actions (gossip, favors, public deeds) tick it. When it fills, they become ``trusted'' −− they can call in a major favor before the timer resets. But enemies also watch their rise, and may work to degrade the timer.

\section{When to Say Yes, No, and Yes, But −− Advanced Rulings}
∈dex{GM rulings!advanced}

At high levels, players will attempt things that break normal bounds. Use these frameworks.

\subsection*{Compromise Solutions}
Instead of a flat ``no'', offer a \textbf{scaled success} with a cost:
\begin{itemize}
  ℹtem \textbf{Yes, but only for this scene.}
  ℹtem \textbf{Yes, but you must mark a permanent Condition.}
  ℹtem \textbf{Yes, but it will take longer −− start a timer.}
  ℹtem \textbf{Yes, but the GM gains 3 SB to spend on that effect.}
\end{itemize}
This keeps creativity alive while respecting the game's tension economy.

\subsection*{Saying ``No'' Without Shutting Down}
If a plan is impossible (e.g., ``I kill the god with a single arrow''), explain why: ``The god is not physically present; you would need to unmake its last shrine first.'' Then offer a path: ``If you find the shrine and perform a ritual, you could banish it −− you'll need to gather three relics.'' Turn the ``no'' into a new arc.

\subsection*{Let the Dice Decide −− But Hedge}
For a truly impossible request, set an \textbf{extreme DV (6+)} and a \textbf{Desperate Position}. If they succeed, they get a lesser version of what they wanted, or a temporary victory. This respects the dice while maintaining verisimilitude.

\begin{tcolorbox}[colback=blue!5!white, colframe=blue!75!black, title=GM Golden Rule for Advanced Play]
\emph{The mechanics are not a ceiling −− they are scaffolding. Build taller stories with them, but never let them crush creativity.}
\end{tcolorbox}